\subsubsection{Translation Invariance as an Inductive Bias}\label{sec:translation-invariance-as-an-inductive-bias}

From the previous section, a convolutional layer has two structural properties: \textbf{sparse connectivity} and \textbf{weight sharing}. The second one, \emph{weight sharing}, has an important consequence: the same feature detector is applied at every spatial location of the input.

\highspace
\begin{flushleft}
    \textcolor{Green3}{\faIcon{clone} \textbf{Reusing the same detector everywhere}}
\end{flushleft}
Suppose a filter $w$ has learned to detect a vertical edge. At one position $(r_1, c_1)$, the convolution computes:
\begin{equation*}
    a(r_1, c_1) = \sum_{u,v,k} w\left(u,v,k\right) \cdot x\left(r_1 + u, c_1 + v, k\right) + b
\end{equation*}
At another position $(r_2, c_2)$, the convolution computes:
\begin{equation*}
    a(r_2, c_2) = \sum_{u,v,k} w\left(u,v,k\right) \cdot x\left(r_2 + u, c_2 + v, k\right) + b
\end{equation*}
The important point is that the parameters $w$ and $b$ are the same in both cases. So if the filter recognizes a useful pattern in the upper-left part of an image, it can recognize the same pattern in the lower-right part as well. Thus, there is \hl{no need to learn a separate set of parameters for every spatial location.}

\highspace
\begin{flushleft}
    \textcolor{Green3}{\faIcon[regular]{lightbulb} \textbf{The underlying assumption of translation invariance}}
\end{flushleft}
The \definition{Translation Invariance} means that \hl{if a feature is useful in one part of the image, it is likely to be useful in other parts as well.} That is the \textbf{inductive bias} imposed by the convolution operation.

\highspace
An \definition{Inductive Bias} is simply an \hl{assumption built into the model architecture about the kind of functions that are likely to be useful.} In this case, the assumption is that \hl{the usefulness of a visual feature should \textbf{not depend strongly} on its absolute position in the image.}

\highspace
For images, this is usually a very sensible prior. For example, a vertical edge is still a vertical edge whether it appears in the upper-left corner of an image or in the lower-right corner. Likewise, a texture or local corner pattern can be meaningful anywhere in the image.

\highspace
\textcolor{Green3}{\faIcon{question-circle} \textbf{Why does weight sharing create this property?}}
Imagine instead a fully connected layer. A neuron looking at pixels in the top-left region would use parameters completely unrelated to a neuron looking at the same pattern in the bottom-right region. Therefore the model could need to learn separately: \emph{pattern at position A}, \emph{pattern at position B}, \emph{pattern at position C}, etc. A convolution forces them all to use the same detector:
\begin{equation*}
    \text{same filter} \to \text{same feature definition everywhere}
\end{equation*}
Hence the spatial position where the pattern appears do not change \textbf{what the filter is looking for}.

\highspace
\textcolor{Green3}{\faIcon{check-circle}} Natural images have a strong spatial regularity, then \hl{the same kinds of local structures can appear almost anywhere in the image.} For example edges, corners, and textures can appear in many different locations. It would therefore be wasteful to learn independent copies of the same detector for every possible image location. \hl{CNNs exploit this prior directly through their architecture.}

\highspace
\textcolor{Green3}{\faIcon{question-circle} \textbf{Why is this called an inductive bias?}} A CNN is not merely learning everything from scratch. Its architecture already tells the network: ``\emph{if a pattern is useful here, try using the same pattern detector elsewhere too}''. So the model search space is deliberately restricted to functions that are likely to be useful for images. Instead of allowing arbitrary weights $W_{\text{arbitrary}}$, we impose a highly structured $W$:
\begin{equation*}
    W_{\text{CNN}} = \text{sparse} + \text{shared}
\end{equation*}
That restriction makes the model less flexible than an arbitrary dense mapping, but in a \textbf{useful direction} for images. This is precisely what an inductive bias does: \hl{restrict the hypothesis space using assumptions about the data.}

\highspace
\begin{takeawaysbox}
    The chain is:
    \begin{equation*}
        \text{weight sharing} \Rightarrow \text{same filter at every spatial location}
    \end{equation*}
    Which introduces the prior that \textbf{a useful visual feature should remain useful when translated}. This is a very appropriate inductive bias for images and simultaneously gives a major parameter-efficiency advantage.

    \highspace
    In other words, CNNs share the same filter weights across spatial locations. This imposes the inductive bias that a feature useful at one image position is also useful elsewhere. As a consequence, the network does not need to learn separate detectors for every position, dramatically reducing the number of parameters.
\end{takeawaysbox}